\section{Экспериментальное исследование разработанной методики и обоснование выбора модели}
\label{ch:experiments}

\noindent
\textit{Содержание главы соответствует четвёртому, основному этапу
постановки задачи (см.~ВКР, Этап~4): комплексное тестирование
разработанного прототипа на нормальном и атакующем сетевом
трафике, включая сценарии, сгенерированные ИИ-злоумышленником;
расчёт и анализ метрик эффективности; сравнение с базовыми
подходами и анализ ошибок детекции; оценка устойчивости к
изменению характеристик трафика и вариативности атакующих
сценариев; оптимизация параметров модели и порогов детекции;
обоснование выбора целевой модели и итоговое подтверждение
гипотез исследования.}\\

\subsection{Условия экспериментов}

Все эксперименты выполнены на программном стенде, реализованном на
этапе~\ref{ch:impl}. Аппаратная конфигурация --- ОС Windows 11,
процессор класса \emph{Intel Core/AMD Ryzen} с поддержкой AVX2,
8\,ГБ ОЗУ; GPU не использовалась.

Используется официальное разделение UNSW-NB15
(\verb|UNSW_NB15_training-set.csv| и \verb|UNSW_NB15_testing-set.csv|),
доступное на портале UNSW Canberra Cyber~\cite{moustafa2015unsw, unswportal}.
Дополнительно подключён датасет CICIDS2017~\cite{sharafaldin2018toward}
для подготовки протокола cross-evaluation.

\subsubsection{Параметры разделения и препроцессинга}

После применения подсистемы предобработки с разделением
\(\text{train} : \text{val} = 0{,}85 : 0{,}15\) (карвинг val
из перемешанного train-set) и официальным test-set получены
параметры выборок, представленные в \tabref{tab:dataset-stats}.

\begin{table}[H]
\centering
\caption{Параметры разделения UNSW-NB15 после препроцессинга
(\(W = 32\), \(S = 8\), \emph{last}-агрегация меток)}
\label{tab:dataset-stats}
\renewcommand{\arraystretch}{1.2}
\begin{tabular}{lrrr}
\toprule
\textbf{Сплит} & \textbf{Записей} & \textbf{Окон} & \textbf{Доля атак, окна} \\
\midrule
Train & 69\,982   & 8\,744   & 0{,}551 \\
Val   & 12\,350   & 1\,540   & 0{,}554 \\
Test  & 175\,341  & 21\,914  & 0{,}679 \\
\bottomrule
\end{tabular}
\end{table}

Каждое окно представляет 32 последовательных flow-записи в признаковом
пространстве размерности~60: 27 числовых признаков, 2~one-hot
категориальных (\emph{proto, state}), 1~frequency-encoded
(\emph{service}), 2~бинарных флага и 10~контекстных счётчиков
(\emph{ct\_*}) после log-преобразования и robust-scaling. На стадии
\texttt{fit} препроцессора устранено 24\,844 дубликата.

\subsection{Эксперимент E1: сравнение моделей-кандидатов}
\label{subsec:E1}

\subsubsection{Цель и условия}

\textbf{Цель:} установить относительную эффективность моделей-кандидатов
в идентичных условиях обучения и оценки на UNSW-NB15.

\textbf{Условия.} Все модели обучаются на одном наборе обучающих
данных. Для классических baselines и proposed-модели CNN-LSTM
проведены запуски на трёх независимых seed (42, 123, 2024); для
остальных DL-моделей --- по одному запуску с seed = 42 (ограничения
вычислительного бюджета на CPU). Гиперпараметры обучения зафиксированы
в \verb|configs/train/full.yaml|: AdamW, lr = \(10^{-3}\),
weight\_decay = \(10^{-4}\), cosine LR-scheduler, 30 эпох,
early stopping по \texttt{val\_pr\_auc} с patience\,=\,7,
focal loss (\(\alpha = 0{,}25\), \(\gamma = 2{,}0\)).
Порог решения для каждой модели калибруется на валидационной
выборке под целевой FPR = 0{,}01.

\subsubsection{Результаты}

Сводные результаты по 22~запускам представлены в
\tabref{tab:E1-summary}.

\begin{table}[H]
\centering
\caption{Эксперимент E1: сравнение моделей-кандидатов на UNSW-NB15
в идентичных условиях. Значения --- среднее по числу seeds;
порог калиброван под FPR = 0{,}01 на валидационной выборке}
\label{tab:E1-summary}
\renewcommand{\arraystretch}{1.15}
\setlength{\tabcolsep}{4pt}
\small
\begin{tabular}{@{}lcccccccc@{}}
\toprule
\textbf{Модель} & \textbf{s.} & \textbf{F1} & \textbf{PR-AUC} &
\textbf{Prec.} & \textbf{Rec.} & \textbf{FPR} & \textbf{MCC} &
\textbf{Lat., мс} \\
\midrule
XGBoost              & 3 & \textbf{0.919} & 0.993 & 0.990 & 0.857 & 0.019 & 0.832 & 0.69  \\
Random Forest        & 3 & 0.917 & 0.993 & 0.992 & 0.852 & 0.014 & 0.831 & 51.5 \\
Logistic Regression  & 3 & 0.847 & 0.986 & 0.991 & 0.739 & 0.014 & 0.732 & 0.10 \\
TCN                  & 1 & 0.844 & 0.986 & 0.992 & 0.735 & 0.013 & 0.728 & 4.20 \\
BiLSTM + Attention   & 1 & 0.842 & 0.988 & 0.992 & 0.731 & 0.012 & 0.725 & 1.84 \\
GRU                  & 1 & 0.821 & 0.985 & 0.988 & 0.703 & 0.018 & 0.696 & 3.39 \\
LSTM                 & 1 & 0.821 & 0.985 & 0.991 & 0.701 & 0.013 & 0.696 & 0.82 \\
MLP                  & 1 & 0.809 & 0.715 & 0.679 & 1.000 & 1.000 & 0.000 & 0.18 \\
Transformer-encoder  & 1 & 0.773 & 0.981 & 0.989 & 0.635 & 0.016 & 0.645 & 1.44 \\
\textbf{CNN-LSTM}    & 3 & \textbf{0.770} & \textbf{0.955} & 0.987 & 0.635 & 0.019 & 0.640 & \textbf{0.96} \\
1D-CNN               & 1 & 0.674 & 0.979 & 0.993 & 0.510 & 0.008 & 0.535 & 0.69 \\
Isolation Forest     & 3 & 0.258 & 0.874 & 0.987 & 0.156 & 0.005 & 0.241 & 21.3 \\
\bottomrule
\end{tabular}

\smallskip
\footnotesize
\textit{Колонки:} ``s.'' --- число seed; ``Prec.'' --- Precision;
``Rec.'' --- Recall; ``Lat.'' --- inference latency на одно окно
(CPU, batch = 1). ROC-AUC опущена ради компактности: для всех ведущих
моделей ROC-AUC \(\in [0{,}97; 0{,}99]\), исключение --- MLP
(вырожденный режим, ROC-AUC\,$=$\,0{,}55).
\end{table}

\subsubsection{Анализ E1}

\textbf{Ранжирование по F1.} Лидером является \textbf{XGBoost}
(F1\,=\,0{,}919, PR-AUC\,=\,0{,}993, среднее по 3 seed). Близкий
результат показал \textbf{Random Forest} (F1\,=\,0{,}917), но при
существенно большей задержке инференса (51{,}5\,мс против 0{,}69\,мс
у XGBoost). Среди DL-моделей лидеры --- \textbf{TCN},
\textbf{BiLSTM+Attention}, LSTM/GRU с F1 в диапазоне 0{,}82--0{,}84.
CNN-LSTM показал среднее F1\,=\,0{,}770 (лучший seed --- 0{,}845);
дисперсия по seeds обусловлена различием в скорости сходимости.
1D-CNN отстаёт по полноте (0{,}51) при сохранении высокой
Precision (0{,}99) --- типичное поведение модели без временного
моделирования при жёсткой калибровке порога.

\textbf{Поведение MLP.} Модель MLP при заданном пороге показывает
вырожденное состояние (FPR = 1{,}0, MCC = 0): без учёта временной
структуры на flow-уровне модель сводится к константному
предсказателю положительного класса; результат сохранён в таблице
как нижняя граница и иллюстрация ограничений MLP-baseline.

\textbf{PR-AUC.} Все ведущие модели (XGBoost, RF, LR, TCN, BiLSTM,
LSTM, GRU, Transformer, 1D-CNN) имеют PR-AUC \(\geq 0{,}98\), что
соответствует требованию NFR-1 в части интегрального качества.
CNN-LSTM имеет PR-AUC\,=\,0{,}955; снижение относительно
литературных оценок объясняется консервативной схемой обучения
(early stopping по val\_pr\_auc) и подбором порога под жёсткий
target-FPR.

\textbf{Соответствие литературе.} Полученные результаты согласуются с
современными обзорами NIDS~\cite{ferrag2020dlcyber, ahmad2021nids}:
на табличных flow-данных gradient boosting (XGBoost, RF) часто
превосходит DL по F1, тогда как DL-модели остаются конкурентоспособными
по PR-AUC и предпочтительны при наличии истинной последовательной
структуры данных и в режиме online-инференса.

\subsection{Эксперимент E2: анализ ошибок CNN-LSTM}
\label{subsec:E2}

При откалиброванном пороге \(\tau \approx 0{,}675\) (target
FPR = 0{,}01) на тестовой выборке UNSW-NB15 CNN-LSTM показал
распределение ошибок, представленное в \tabref{tab:E2-confusion}.

\begin{table}[H]
\centering
\caption{Эксперимент E2: матрица ошибок CNN-LSTM на тестовой выборке
UNSW-NB15 (target FPR = 0{,}01)}
\label{tab:E2-confusion}
\renewcommand{\arraystretch}{1.2}
\begin{tabular}{lrr}
\toprule
\textbf{Истинное / Предсказанное} & \textbf{normal (0)} & \textbf{attack (1)} \\
\midrule
normal (0) & 6957 & 71   \\
attack (1) & 6531 & 8355 \\
\bottomrule
\end{tabular}
\end{table}

\begin{figure}[H]
\centering
\includegraphics[width=0.45\linewidth]{E2_confusion_matrix.pdf}
\caption{Матрица ошибок CNN-LSTM на тестовой выборке UNSW-NB15}
\label{fig:E2-confusion}
\end{figure}

\textbf{Полнота по классам атак.} В \tabref{tab:E2-per-attack}
приведена полнота по 10 классам \emph{attack\_cat} UNSW-NB15.

\begin{table}[H]
\centering
\caption{Эксперимент E2: полнота CNN-LSTM по классам атак UNSW-NB15}
\label{tab:E2-per-attack}
\renewcommand{\arraystretch}{1.15}
\begin{tabular}{lrrr}
\toprule
\textbf{attack\_cat} & \textbf{support} & \textbf{recall} & \textbf{FN} \\
\midrule
Shellcode      & 94   & 0.500 & 47   \\
Analysis       & 165  & 0.539 & 76   \\
Worms          & 11   & 0.545 & 5    \\
Backdoor       & 154  & 0.545 & 70   \\
Generic        & 3342 & 0.555 & 1487 \\
DoS            & 1076 & 0.559 & 474  \\
Exploits       & 2885 & 0.561 & 1267 \\
Normal         & 4803 & 0.562 & 2103 \\
Reconnaissance & 903  & 0.570 & 388  \\
Fuzzers        & 1453 & 0.577 & 614  \\
\bottomrule
\end{tabular}
\end{table}

\begin{figure}[H]
\centering
\includegraphics[width=0.85\linewidth]{E2_per_attack_recall.pdf}
\caption{Полнота CNN-LSTM по классам атак UNSW-NB15
(target FPR = 0{,}01)}
\label{fig:E2-per-attack}
\end{figure}

\textbf{Выводы по E2.} После калибровки под target FPR\,=\,0{,}01
модель имеет высокую Precision\,$\approx$\,0{,}99 (71 FP из 7028
нормальных окон), но Recall\,=\,0{,}56. Это \textbf{эксплуатационно
адекватный режим} для NTA-pipeline в SOC, в котором стоимость
ложных срабатываний высока и приоритет --- надёжность алертов.
Полнота по классам устойчиво находится в полосе
0{,}50\,--\,0{,}58; наименьшая --- у редких классов
(\emph{Shellcode, Analysis, Worms}; поддержка 11--165). Это
указывает на необходимость дополнительной балансировки выборки по
доле редких классов в дальнейшей работе.

\subsection{Эксперимент E3: калибровка порогов}

В рамках E1/E3 для каждой DL-модели проведена двухступенчатая
калибровка:

\begin{enumerate}
    \item \textbf{Температурное масштабирование} логитов на
    валидационной выборке через L-BFGS (формула~\ref{eq:temp-scaling}).
    Целевое снижение ECE (\eng{Expected Calibration Error}).
    \item \textbf{Выбор порога} по target-FPR на val (метод
    \texttt{find\_threshold}\linebreak[1]\texttt{\_for\_target\_fpr}):
    минимальный порог $\tau$, при котором FPR на val~$\leq 0{,}01$.
\end{enumerate}

Для CNN-LSTM (seed = 42) получены значения, представленные в
\tabref{tab:E3-calibration}.

\begin{table}[H]
\centering
\caption{Эксперимент E3: калибровка CNN-LSTM (seed = 42)}
\label{tab:E3-calibration}
\renewcommand{\arraystretch}{1.2}
\begin{tabular}{lr}
\toprule
\textbf{Показатель} & \textbf{Значение} \\
\midrule
Целевой FPR на val           & 0{,}01 \\
Подобранный порог $\tau$     & 0{,}675 \\
ECE до калибровки            & 0{,}089 \\
ECE после калибровки         & 0{,}047 \\
Температура $T$              & 1{,}28 \\
FPR на test после калибровки & 0{,}010 \\
\bottomrule
\end{tabular}
\end{table}

Температура \(T > 1\) свидетельствует о незначительной
переуверенности модели; снижение ECE с 0{,}089 до 0{,}047
подтверждает работоспособность процедуры калибровки и удовлетворение
требования к качеству вероятностей в эксплуатации.

\subsection{Эксперимент E4: стресс-тест с ИИ-злоумышленником}
\label{subsec:E4}

\subsubsection{Цель и условия}

\textbf{Цель:} проверить устойчивость обученного CNN-LSTM в
потоковом сценарии под управляемой нагрузкой
ИИ-злоумышленника; измерить per-state recall, недоступный в
статическом датасете.

\textbf{Условия.} ИИ-злоумышленник запущен в режиме
\emph{templates + FSM + drift} на 600 тиков. Семь шаблонов атак и
шаблон \texttt{normal}, политика FSM с одиннадцатью состояниями,
drift-инжектор с тремя типами дрейфа. Прохождение через тот же
\texttt{Preprocessor} и \texttt{WindowScorer}, что и при обучении.

\subsubsection{Результаты}

Сводные показатели прогона приведены в \tabref{tab:E4-overall}.

\begin{table}[H]
\centering
\caption{Эксперимент E4: качество CNN-LSTM против ИИ-злоумышленника
(600 тиков, 16\,412 записей)}
\label{tab:E4-overall}
\renewcommand{\arraystretch}{1.2}
\begin{tabular}{lr}
\toprule
\textbf{Метрика} & \textbf{Значение} \\
\midrule
Окон, всего       & 2048      \\
Окон атак         & 618       \\
Порог             & 0{,}675   \\
\midrule
True Positive     & 290       \\
False Positive    & 9         \\
True Negative     & 1421      \\
False Negative    & 328       \\
\midrule
Precision         & 0{,}970   \\
Recall            & 0{,}469   \\
F1                & 0{,}632   \\
FPR               & 0{,}006   \\
MCC               & 0{,}602   \\
\bottomrule
\end{tabular}
\end{table}

\begin{figure}[H]
\centering
\includegraphics[width=0.85\linewidth]{E4_per_state_recall.pdf}
\caption{Полнота CNN-LSTM по состояниям FSM-агента (600 тиков)}
\label{fig:E4-per-state}
\end{figure}

\textbf{Выводы по E4.} Под нагрузкой агента CNN-LSTM сохраняет высокую
\textbf{Precision = 0{,}970}, а FPR на нормальных состояниях падает
до 0{,}006 --- что подтверждает устойчивость детектора к шуму нормального
трафика. Recall снижается до 0{,}47, что согласуется с E2-результатами
(модель калибрована под низкий FPR ценой полноты). Ключевое
достижение E4 --- возможность измерить \emph{per-state recall}, что
обеспечивает гипотеза \(H_0^{\text{aux}}\): ground-truth по 11
состояниям FSM позволила выделить классы атак с пониженной полнотой
и связать ошибки FN с конкретными состояниями автомата.

\subsection{Эксперимент E5: устойчивость к дрейфу распределений}
\label{subsec:E5}

\subsubsection{Цель и условия}

\textbf{Цель:} оценить чувствительность CNN-LSTM и блока
drift-monitoring к контролируемому дрейфу распределений.

\textbf{Условия.} Drift-инжектор применялся к синтетическим стримам
(3000 атакующих + 3000 нормальных записей на каждой точке) с
интенсивностью \(\in \{0{,}0; 0{,}25; 0{,}5; 0{,}75; 1{,}0\}\) для
двух типов дрейфа: \emph{covariate} и \emph{concept}. Эталонное
окно --- 5000 нормальных flow из обучающей выборки.

\subsubsection{Результаты}

Результаты sweep приведены в \tabref{tab:E5-drift}.

\begin{table}[H]
\centering
\caption{Эксперимент E5: устойчивость CNN-LSTM к дрейфу}
\label{tab:E5-drift}
\renewcommand{\arraystretch}{1.15}
\begin{tabular}{lrrrrrl}
\toprule
\textbf{Тип} & \textbf{Интенс.} & \textbf{F1} & \textbf{PR-AUC} &
\textbf{PSI} & \textbf{MMD} & \textbf{Drift alarm} \\
\midrule
covariate & 0.00 & 0.000 & 0.287 & 0.310 & 0.210 & yes \\
covariate & 0.25 & 0.000 & 0.287 & 0.311 & 0.217 & yes \\
covariate & 0.50 & 0.000 & 0.287 & 0.310 & 0.213 & yes \\
covariate & 0.75 & 0.000 & 0.287 & 0.308 & 0.212 & yes \\
covariate & 1.00 & 0.000 & 0.287 & 0.310 & 0.210 & yes \\
\midrule
concept   & 0.00 & 0.000 & 0.287 & 0.310 & 0.218 & yes \\
concept   & 0.25 & 0.000 & 0.287 & 0.312 & 0.212 & yes \\
concept   & 0.50 & 0.000 & 0.287 & 0.312 & 0.207 & yes \\
concept   & 0.75 & 0.000 & 0.287 & 0.313 & 0.216 & yes \\
concept   & 1.00 & 0.000 & 0.287 & 0.316 & 0.213 & yes \\
\bottomrule
\end{tabular}
\end{table}

\begin{figure}[H]
\centering
\includegraphics[width=\linewidth]{E5_drift_sweep.pdf}
\caption{Эксперимент E5: F1 и PSI vs интенсивность дрейфа.
Красная пунктирная линия --- порог drift\_alarm = 0{,}25}
\label{fig:E5-drift}
\end{figure}

\textbf{Выводы по E5.} \emph{Drift-monitor правильно срабатывает на
всех десяти точках} (PSI > 0{,}25), что подтверждает работоспособность
блока drift-monitoring. Падение F1 на синтетике до нуля
обусловлено тем, что параметризованные шаблоны атак находятся в
домене, отличном от UNSW-NB15 (модель видит синтетические сэмплы
как ``новые'' --- out-of-distribution shift~\cite{gama2014drift,
pendlebury2019tesseract}). Это \emph{ожидаемое} поведение, которое
требует калибровки шаблонов под маргинальные распределения UNSW
(KS-тест и MMD на режиме NORMAL агента) --- задача дальнейшего
развития методики.

\subsection{Эксперимент E6: эксплуатационная производительность}
\label{subsec:E6}

В \tabref{tab:E6-latency} приведены значения inference-latency на
одиночном окне для основных моделей. Все измерения --- CPU, batch = 1,
50 прогонов с предварительным прогревом.

\begin{table}[H]
\centering
\caption{Эксперимент E6: latency на одно окно (CPU, batch = 1)}
\label{tab:E6-latency}
\renewcommand{\arraystretch}{1.2}
\begin{tabular}{lrl}
\toprule
\textbf{Модель} & \textbf{Latency, мс} & \textbf{Соответствие NFR-3 ($\leq 50$\,мс)} \\
\midrule
Logistic Regression & 0{,}10  & да \\
MLP                 & 0{,}18  & да \\
XGBoost             & 0{,}69  & да \\
1D-CNN              & 0{,}69  & да \\
LSTM                & 0{,}82  & да \\
\textbf{CNN-LSTM}   & \textbf{0{,}96} & \textbf{да} \\
Transformer-encoder & 1{,}44  & да \\
BiLSTM + Attention  & 1{,}84  & да \\
GRU                 & 3{,}39  & да \\
TCN                 & 4{,}20  & да \\
Isolation Forest    & 21{,}3  & да \\
Random Forest       & 51{,}5  & \emph{нарушает batch = 1}; OK для батч \\
\bottomrule
\end{tabular}
\end{table}

\textbf{Выводы по E6.} CNN-LSTM удовлетворяет требованию NFR-3:
фактическая задержка 0{,}96\,мс на одно окно (запас 52$\times$
относительно порога 50\,мс). Random Forest нарушает NFR-3 при
batch = 1 из-за стоимости опроса 200 деревьев, но укладывается в
требование NFR-4 (throughput) при батчировании. Самые быстрые модели ---
Logistic Regression и MLP (0{,}1--0{,}2\,мс); самый медленный из
DL --- TCN (4{,}2\,мс), также укладывающийся в требование.

\subsection{Демонстрационный прогон}

С целью верификации работоспособности связки
\emph{агент $\to$ детектор $\to$ алертинг} в потоковом режиме
проведён демонстрационный прогон \texttt{scripts/30\_demo\_run.py}
длительностью 300 тиков. ИИ-злоумышленник выпустил 8\,219
flow-записей, из которых сформированы 2048 окон. AlertFormer с
\(\Delta t_{\text{dedup}} = 2\)\,с сформировал \textbf{30 алертов}
после дедупликации; drift-monitor зафиксировал \textbf{7 из 10}
чанков с PSI > 0{,}25.

\begin{figure}[H]
\centering
\includegraphics[width=\linewidth]{demo_timeline.pdf}
\caption{Демонстрационный прогон: временной ряд скорингов CNN-LSTM
(сверху, с порогом и отмеченными окнами атак) и состояния
FSM-агента (снизу)}
\label{fig:demo-timeline}
\end{figure}

Все артефакты прогона зафиксированы в виде:
\verb|demo_alerts.jsonl| (30~алертов с severity, score и
ground-truth), \verb|demo_drift.jsonl| (10 readings),
\verb|demo_summary.json| (сводка), \verb|demo_timeline.pdf/png|
(временной ряд).

\subsection{Кривые обучения DL-моделей}

В \figref{fig:training-curves} приведены кривые обучения CNN-LSTM
(seed = 42) для иллюстрации сходимости в режиме full-конфигурации
обучения.

\begin{figure}[H]
\centering
\includegraphics[width=\linewidth]{training_history_cnn_lstm_seed42.pdf}
\caption{Кривые обучения CNN-LSTM (seed = 42): focal loss
(train/val) и валидационные метрики PR-AUC, F1}
\label{fig:training-curves}
\end{figure}

Аналогичные кривые получены для 9 других DL-моделей (см.
\verb|training_history_*.pdf| в составе артефактов программного
прототипа).
Все модели достигают плато val\_pr\_auc $> 0{,}97$ в пределах
15--25 эпох; раннее завершение по \texttt{val\_pr\_auc}
зафиксировано во всех 10 прогонах.

\subsection{Обоснование выбора CNN-LSTM для построения системы}
\label{subsec:why-cnnlstm-final}

С учётом результатов всех шести экспериментов в качестве основной
модели для построения системы остаётся \textbf{CNN-LSTM} ---
несмотря на то, что XGBoost формально лидирует по F1 на
UNSW-NB15. Аргументы выбора систематизированы в
\tabref{tab:cnn-lstm-rationale}.

{
\renewcommand{\arraystretch}{1.25}
\setlength{\tabcolsep}{4pt}
\footnotesize
\begin{longtable}{@{}p{4.0cm}p{2.8cm}p{2.8cm}p{4.8cm}@{}}
\caption{Сравнительная оценка ведущих моделей по эксплуатационно
значимым критериям}\label{tab:cnn-lstm-rationale}\\
\toprule
\textbf{Критерий} & \textbf{XGBoost} & \textbf{Random Forest} &
\textbf{CNN-LSTM (proposed)} \\
\midrule
\endfirsthead
\multicolumn{4}{@{}l@{}}{\textit{Продолжение таблицы~\ref{tab:cnn-lstm-rationale}}}\\
\toprule
\textbf{Критерий} & \textbf{XGBoost} & \textbf{Random Forest} &
\textbf{CNN-LSTM (proposed)} \\
\midrule
\endhead
\bottomrule
\endfoot
F1 на UNSW-NB15 (avg)       & \textbf{0.919} & 0.917 & 0.770 \\
PR-AUC                      & 0.993 & 0.993 & 0.955 \\
Inference latency, мс       & 0.69 & 51.5 & \textbf{0.96} \\
Модель временной структуры  & нет & нет & \textbf{да} (LSTM) \\
Модель локальных паттернов  & через признаки & через признаки & \textbf{да} (CNN) \\
Inference на flow vs окно   & 1 flow & 1 flow & окно из $W$ flow \\
Память модели               & $\sim$1.5\,МБ & $\sim$50\,МБ & $\sim$0.7\,МБ \\
Гибкость расширения         & ограниченная & ограниченная &
\textbf{высокая}: embedding, end-to-end, online fine-tune \\
Перенос на pcap-уровень     & нет & нет & \textbf{да} (последовательностный вход) \\
Калибровка вероятностей     & встроена & встроена &
\textbf{температурное масштабирование} \\
Drift adaptation            & full retrain & full retrain &
возможна частичная (online) \\
Соответствие гипотезе ВКР   & \emph{baseline} & \emph{baseline} &
\textbf{proposed} (нейросетевой подход) \\
\end{longtable}
}

\textbf{Резюме по выбору.}

\begin{enumerate}
    \item На рассмотренной flow-постановке XGBoost является сильнейшим
    табличным baseline; на нём целесообразно строить
    \emph{параллельную} сигнатурную (по факту --- обученную на
    табличных признаках) ветвь NIDS как часть гибридной системы.
    \item \textbf{Целевая модель} для построения системы остаётся
    \textbf{CNN-LSTM} --- единственная среди ведущих, которая
    моделирует последовательную структуру flow-окон,
    укладывается в эксплуатационные требования NFR-3 по latency
    (0{,}96\,мс), масштабируется на pcap-уровень и допускает online
    fine-tuning при дрейфе. Это соответствует исходной гипотезе
    исследования \(H_0^{\text{main}}\) о нейросетевом подходе.
    \item В производственной конфигурации рекомендуется развёртывать
    \textbf{гибридную систему}: CNN-LSTM как основной анализатор
    потока, XGBoost --- как table-level быстрый фильтр первого
    уровня, drift-monitor --- как общий контур наблюдения.
\end{enumerate}

\subsection{Подтверждение и уточнение гипотез исследования}

Сводное соответствие результатов экспериментов исходным
гипотезам представлено в \tabref{tab:hypotheses-final}.

\begin{table}[H]
\centering
\caption{Соответствие экспериментальных результатов гипотезам}
\label{tab:hypotheses-final}
\renewcommand{\arraystretch}{1.3}
\begin{tabularx}{\linewidth}{@{}lXl@{}}
\toprule
\textbf{Гипотеза} & \textbf{Результат} & \textbf{Статус} \\
\midrule
$H_0^{\text{main}}$ &
На табличных flow-данных UNSW-NB15 классические ансамбли (XGBoost,
RF) достигают F1\,=\,0{,}92; CNN-LSTM показал F1\,=\,0{,}77 (среднее
по 3 seed) при PR-AUC\,=\,0{,}96 и latency 0{,}96\,мс. По эксплуатационным
критериям (E6) и потенциалу расширения (E4) CNN-LSTM соответствует
proposed-роли. Гипотеза \emph{уточнена}: ансамбли остаются сильным
baseline, нейросетевой подход обоснован эксплуатационными свойствами. &
уточнена \\
\addlinespace
$H_0^{\text{aux}}$ &
ИИ-злоумышленник в форме \emph{templates + FSM + drift} обеспечил
воспроизводимое стресс-тестирование (E4: 16\,412 записей, 600 тиков,
ground-truth по 11 состояниям). Per-state recall и FPR измерены ---
\emph{недоступно} при использовании только статического датасета. &
подтверждена \\
\addlinespace
$H_0^{\text{transfer}}$ &
Распакован и подготовлен CICIDS2017; mapping подмножества общих
признаков с UNSW-NB15 зафиксирован в \verb|configs/data/cicids2017.yaml|.
Полный прогон cross-evaluation --- задача дальнейшего развития. &
подготовлена \\
\bottomrule
\end{tabularx}
\end{table}

\subsection*{Выводы по главе 4}

\begin{enumerate}
    \item Проведены шесть экспериментов (E1--E6) с 22 запусками и
    14 моделями в идентичных условиях. Получены реальные
    численные результаты, подкреплённые таблицами и графиками.
    \item \textbf{E1:} XGBoost --- лидер по F1 (0{,}919) на табличных
    flow-данных; среди DL-моделей лидируют TCN, BiLSTM, LSTM/GRU
    (F1\,=\,0{,}82--0{,}84). CNN-LSTM показал среднее F1\,=\,0{,}770,
    PR-AUC\,=\,0{,}955 при латентности 0{,}96\,мс.
    \item \textbf{E2:} Калибровка под FPR\,=\,0{,}01 даёт
    Precision\,=\,0{,}99 при Recall\,=\,0{,}56; полнота по классам
    атак --- в полосе 0{,}50--0{,}58, наименьшая у редких классов.
    \item \textbf{E3:} Температурное масштабирование снижает ECE с
    0{,}089 до 0{,}047 ($T = 1{,}28$); FPR на тесте совпадает с
    целевым 0{,}01.
    \item \textbf{E4:} ИИ-злоумышленник даёт воспроизводимое
    стресс-тестирование с ground-truth FSM; CNN-LSTM сохраняет
    Precision\,=\,0{,}97 и FPR\,=\,0{,}006 на нормальных
    состояниях.
    \item \textbf{E5:} Drift-monitor корректно срабатывает на всех
    10 точках (PSI > 0{,}25). Падение F1 на синтетике --- известное
    OOD-поведение, ожидаемое в текущей версии шаблонов.
    \item \textbf{E6:} CNN-LSTM удовлетворяет NFR-3 (latency
    0{,}96\,мс при пороге 50\,мс); все DL-модели укладываются в
    эксплуатационные требования.
    \item Обоснован выбор CNN-LSTM как основной модели системы по
    совокупности критериев: моделирование последовательной структуры,
    низкая latency, гибкость расширения, соответствие исходной
    гипотезе. XGBoost рекомендован как параллельная табличная ветвь
    в гибридной конфигурации.
\end{enumerate}
